\section{引言}

智能体帮助用户取消一笔订单时，需要向订单系统传入订单编号。但如果它同时把用户的出生日期和支付卡号也填进了可选参数，订单一样能取消成功，敏感信息却已经发给了本不需要它的服务。这种披露既不是攻击者诱导的，也不影响任务完成——它只是智能体在拥有完整用户档案时，多传了几个字段。

本文将这种现象称为\textbf{冗余参数泄露}（Redundant Parameter Leakage, RPL）。它与提示注入攻击不同\citep{debenedetti2024agentdojo}：后者是外部攻击者试图改变智能体行为，前者发生在正常的、功能正确的任务执行中。它也不是参数填错——工具选对了、必填参数齐全、最终状态也正确，问题只是调用里多了几个本可删掉的敏感参数。

隐私研究早已指出，信息是否可以披露不能只看字段是否敏感，还要看它流向谁、为了什么目的\citep{nissenbaum2004contextual}。同一个电话号码，通知工具可能需要，但查询工具不需要。CONFAIDE 和 PrivacyLens 等工作进一步发现，模型嘴上说知道隐私规范，手上却可能照做不误\citep{mireshghallah2024confaide,shao2024privacylens}。\ToolPrivacyBench 把这一思路带进工具调用场景，发现任务完成率和隐私合规并不是一回事\citep{hu2026toolprivacybench}。AgentSCOPE 追踪信息在用户、智能体和工具之间的传播路径\citep{ngong2026agentscope}，Privacy in Action 关注动态工具生态中的判断与行动差距\citep{wang2025privacyinaction}。这些工作说明，只看最终回答或任务是否完成，不足以发现执行过程中的参数级泄露。

但要把现有工具调用任务改造成能评测 RPL 的数据，并不容易。$\tau$-bench、BFCL、API-Bank 和 AppWorld 四个基准各自保存参考调用的方式不同：有的是任务动作，有的是多轮对话，有的是逐条 API 记录，有的是通过环境验证的执行日志。它们对用户身份、工具可见范围和轨迹完整性的表达也不一样。如果转换时直接把这些记录都叫"标准答案"，不仅会丢掉来源信息，还可能把一条参考路径误当成唯一解或最短路径。

本文的做法是：不改变上游参考调用的任何参数，只在此基础上构造一个"多传了敏感字段"的对照版本。如图\ref{fig:rpl-motivation}所示，两条调用使用同一工具、同一必填参数，唯一区别是后者增加了若干可选的敏感字段。这样，隐私泄露就与工具选择错误、必填参数缺失等功能问题分离开来，差异精确落在多出来的参数上。

\begin{figure}[t]
  \centering
  \begin{tikzpicture}[
    font=\small,
    card/.style={rounded corners=2pt, draw, thick, align=left, inner sep=5pt},
    field/.style={font=\footnotesize, anchor=west},
    hdr/.style={font=\small\bfseries, anchor=west},
    arrow/.style={-{Latex[length=2mm]}, thick},
  ]
    % ---- 用户档案 ----
    \node[card, draw=RPLBlue, fill=RPLBlue!5, text width=4.2cm] (profile) {};
    \node[hdr] at ([xshift=5pt,yshift=-5pt]profile.north west) {用户档案};
    \node[field] at ([xshift=5pt,yshift=-20pt]profile.north west) {user\_id: sophia\_silva};
    \node[field] at ([xshift=5pt,yshift=-32pt]profile.north west) {email: s.silva@mail.com};
    \node[field] at ([xshift=5pt,yshift=-44pt]profile.north west) {date\_of\_birth: 1996-05-12};
    \node[field] at ([xshift=5pt,yshift=-56pt]profile.north west) {payment\_method\_id: cc\_7334};
    \node[field] at ([xshift=5pt,yshift=-68pt]profile.north west) {home\_address: 12 Oak St};

    % ---- 任务指令 ----
    \node[card, draw=RPLDark, fill=gray!6, text width=4.2cm, below=12mm of profile] (task) {};
    \node[hdr] at ([xshift=5pt,yshift=-5pt]task.north west) {任务};
    \node[field, text width=3.8cm] at ([xshift=5pt,yshift=-20pt]task.north west) {取消航班预订 IFOYYZ};

    % ---- Ground call ----
    \node[card, draw=RPLGreen, fill=RPLGreen!5, text width=4.6cm, below=10mm of task, xshift=-3.1cm] (ground) {};
    \node[hdr, text=RPLGreen] at ([xshift=5pt,yshift=-5pt]ground.north west) {不泄露调用};
    \node[field] at ([xshift=5pt,yshift=-20pt]ground.north west) {tool: cancel\_reservation};
    \node[field] at ([xshift=5pt,yshift=-32pt]ground.north west) {reservation\_id: IFOYYZ};

    % ---- Leaking call ----
    \node[card, draw=RPLRed, fill=RPLRed!4, text width=4.6cm, below=10mm of task, xshift=3.1cm] (leak) {};
    \node[hdr, text=RPLRed] at ([xshift=5pt,yshift=-5pt]leak.north west) {泄露调用};
    \node[field] at ([xshift=5pt,yshift=-20pt]leak.north west) {tool: cancel\_reservation};
    \node[field] at ([xshift=5pt,yshift=-32pt]leak.north west) {reservation\_id: IFOYYZ};
    \node[field, text=RPLRed] at ([xshift=5pt,yshift=-44pt]leak.north west) {+ date\_of\_birth: 1996-05-12};
    \node[field, text=RPLRed] at ([xshift=5pt,yshift=-56pt]leak.north west) {+ payment\_method\_id: cc\_7334};

    % ---- 结果 ----
    \node[card, draw=gray, fill=gray!4, text width=3.2cm, below=8mm of ground] (ok1) {};
    \node[hdr] at ([xshift=5pt,yshift=-5pt]ok1.north west) {结果};
    \node[field] at ([xshift=5pt,yshift=-20pt]ok1.north west) {预订已取消};

    \node[card, draw=gray, fill=gray!4, text width=3.2cm, below=8mm of leak] (ok2) {};
    \node[hdr] at ([xshift=5pt,yshift=-5pt]ok2.north west) {结果};
    \node[field] at ([xshift=5pt,yshift=-20pt]ok2.north west) {预订已取消};

    % ---- 箭头 ----
    \draw[arrow, RPLBlue] (profile) -- (task);
    \draw[arrow, RPLGreen] (task) -- node[left, font=\footnotesize] {必要参数} (ground);
    \draw[arrow, RPLRed] (task) -- node[right, font=\footnotesize] {+ 冗余参数} (leak);
    \draw[arrow, gray] (ground) -- (ok1);
    \draw[arrow, gray] (leak) -- (ok2);

    % ---- 标注 ----
    \node[font=\footnotesize\itshape, text=RPLRed, below=2mm of ok2, xshift=-1.55cm, text width=7cm, align=center]
      {RPL\_trigger\_params = \{date\_of\_birth, payment\_method\_id\}\\\textup{= keys(leaking) $-$ keys(ground)}};
  \end{tikzpicture}
  \caption{RPL 配对调用示例。不泄露调用只包含取消预订所需的 \texttt{reservation\_id}；泄露调用保留该参数不变，额外增加两个可选敏感字段。两次调用都成功取消预订，但后者向预订系统披露了不必要的出生日期和支付方式。触发参数精确等于两次调用的参数键差集。}
  \label{fig:rpl-motivation}
\end{figure}

围绕这一目标，本文回答三个问题：

\begin{enumerate}[label=\textbf{RQ\arabic*:}]
  \item 如何把四个结构各异的基准统一成同一套可审计的任务、工具、用户档案和调用表示，同时不丢失来源差异？
  \item 怎样保证每个泄露参数确实来自当前用户、在工具中是可选的、并且不影响原有调用？
  \item 如何区分"调用和源数据一致""调用能跑通""最终状态正确"这三种不同程度的验证，避免把一种证据当成完整证明？
\end{enumerate}

本文的主要贡献如下：

\begin{enumerate}
  \item 提出冗余参数泄露的配对构造方法：每步提供不泄露和泄露两个调用版本，差异只落在多出的可选敏感参数上，可以精确计算。
  \item 分析 $\tau$-bench、BFCL、API-Bank 与 AppWorld 的原始任务结构和验证机制，设计保留来源差异的适配与审计流程。
  \item 构建包含 328 条任务的 RPL 数据，逐条提供工具定义、用户档案、配对调用和来源证据。
  \item 分层报告结构校验、源重放、环境执行和最终状态验证的结果，并明确哪些是已验证的、哪些仍不是定论。
\end{enumerate}

本文不宣称已经获得人工业务授权意义上的最终标准答案，而是提供一个透明、保守、可重复运行的参数级隐私评测基础。下文依次介绍相关研究、源基准、数据构造、验证结果与局限。
